% =========================================================
\section{Related Work}
\label{sec:related_work}
% =========================================================

This section reviews Koopman modeling and recursive prediction, temporal sequence modeling, and real-time predictive control, and positions the proposed method.

\subsection{Koopman Learning and Recursive Prediction}

Koopman operator theory represents nonlinear dynamics as a linear evolution on the observable functions of the state \cite{koopman1931hamiltonian, koopman1932spectra, mezic2005spectral}, which enables the direct application of standard linear control synthesis to nonlinear systems. Early extensions such as DMDc \cite{proctor2016dmdc}, input-aware Koopman formulations \cite{proctor2018generalizing, brunton2016invariant}, and Koopman-based MPC \cite{korda2018linear} established a foundation for applying Koopman theory in control, and later work added learned representations and performance guarantees \cite{han2020deep, shi2022deep, strasser2026overview, duran2025control}. Across this literature, controller performance is highly sensitive to the quality of the underlying model, which calls for improved lifting design.

To learn Koopman models from trajectory data, EDMD fits observables from a pre-defined dictionary such as polynomials or Fourier bases \cite{williams2015data, williams2015kernel, klus2018data}, whereas data-driven methods replace hand-crafted dictionaries with learned representations \cite{takeishi2017koopman, lusch2018deep, yeung2019learning, morton2019deep, yuan2022variational}. Most of these methods construct observables from instantaneous states and identify a single time-invariant operator, neglecting short-term temporal correlations and limiting prediction accuracy over long horizons.

Two lines of work relax this restriction. First, delay-embedding variants enrich the observable dictionary with time-shifted copies of the measurements, either as explicit time-delay observables \cite{kamb2020timedelay} or through Hankel/HAVOK coordinates \cite{brunton2017chaos, arbabi2017ergodic, das2019delay}; the identified operator then acts on a history-augmented state, but history enters only through the observables. Second, context-aware and time-varying formulations let the effective operator itself depend on recent data: Koopa computes context-aware operators within a temporal neighborhood for non-stationary forecasting \cite{liu2023koopa}, and MamKO generates time-varying Koopman matrices from online history for predictive control \cite{li2025mamko}. Our formulation also lets history influence the transition, but confines it to a single low-rank, structure-parameterized family identified jointly with the lifting map and linear in the lifted state, so that the predictive-control problem remains a quadratic program.

Multi-step training has also been adopted in learned Koopman models to mitigate error accumulation under recursive prediction \cite{lusch2018deep, otto2019linearly, azencot2020forecasting, wu2026multistep}; multi-step simulation-error objectives are likewise standard in neural system identification more broadly \cite{forgione2021dynonet}. We therefore do not regard multi-step prediction loss itself as a methodological novelty. The issue considered here is more specific to history-conditioned dynamics: during recursive deployment the effective lifted transition is repeatedly rescheduled according to the predicted history, so both the propagated latent state and the conditioning context can depart from the encoder-consistent one-step training distribution. Our rollout-aware procedure addresses this training--deployment inconsistency by reproducing the recursively scheduled propagation during training, rather than merely extending a one-step objective over multiple prediction horizons.

\subsection{Transformers for Temporal Sequence Modeling}

Transformers \cite{vaswani2017attention} model temporal dependencies via self-attention, capturing pairwise dependencies across all time steps simultaneously and allowing more parallelization and better scalability than step-by-step recurrent networks. Efficient variants and patch-based tokenization \cite{zhou2021informer, wang2020linformer, choromanski2021performer, wu2021autoformer, zhou2022fedformer, nie2023time} further improve efficiency by aggregating local temporal structure before global attention. In the context of Koopman lifting, where the relevant temporal context is not known a priori, Transformers can incorporate short-term temporal information from past trajectories.

Recent work has explored Transformer encoders with Koopman-style latent dynamics for forecasting \cite{forootani2026deepkoopformer, forootani2026spectral}. In particular, \cite{forootani2026spectral} parameterizes the transition as a single low-rank operator with a spectrally constrained diagonal factor on a patch-based Transformer backbone. This line shares our use of a patch-based Transformer encoder and of a low-rank factorization, but differs in two respects that matter for control: the operator there is a single static low-rank transition used for forecasting, whereas we keep a full-rank nominal transition and let history modulate only a low-rank deviation; and our prediction is generated by the identified Koopman model inside a real-time predictive controller rather than by the Transformer itself. Relatedly, low-rank space-varying transitions have been learned for continuum-robot control, with auxiliary networks parameterizing the state dependence of the transition spectrum and of the input matrix \cite{feizi2025deepkoopman}; that formulation is conditioned on the instantaneous state, modulates both the transition and the input matrix, and is paired with an infinite-horizon linear-quadratic controller rather than a constrained predictive one.

\subsection{Real-Time and Scheduled Predictive Control}

Two established strategies reduce the online cost of predictive control. The real-time iteration (RTI) performs a single sequential quadratic programming step per sampling instant around the current nominal trajectory, so that a quadratic program replaces the full nonlinear program \cite{diehl2005realtime, gros2020rti}; its nominal stability and the effect of suboptimality have been analyzed \cite{diehl2005rtistability}, and sensitivity-based schemes reduce the computational delay by separating an offline nonlinear program from a fast online update \cite{zavala2009advancedstep}. Scheduled, or quasi-linear parameter-varying (quasi-LPV), predictive control instead embeds a nonlinear model into a family of linear models indexed by a scheduling variable, and freezing or iterating the scheduling trajectory turns the nonlinear problem into a sequence of QPs \cite{rugh2000gain, leith1998velocity, cisneros2016qlmpc, cisneros2020qlpv}.

Our controller inherits the one-iteration RTI structure but differs in where the scheduling signal comes from. In quasi-LPV MPC the scheduling variable is an explicit function of the state or of the velocity, and its predicted trajectory is obtained from the previous solution or from a planned reference. Here the effective transition is scheduled by a context that a learned encoder extracts from a history window, and during prediction the history is itself produced by the predictor, so the scheduling trajectory is endogenous to the learned model and cannot be read off the measured state. Two consequences are exploited in Section~\ref{sec:control}: the scheduling context must be re-evaluated along the rolled-out latent trajectory rather than assumed, and the full encoder need not be re-run at every prediction stage, which allows the single-iteration controller to fit inside the sampling period.

